\noindent\textbf{Abstract}\par\smallskip
\noindent Force-plate software summarises each countermovement jump (CMJ) with a few vendor metrics, and athletes with similar values are treated as equivalent. We tested that assumption with \nTrials{} routine CMJ trials from \nAth{} male footballers in India. Principal component analysis of body-weight-normalised take-off and landing force--time curves retained four components (\cumFour\% of variance), concentrated in the landing phase. Trials matched on jump height or relative peak force varied nearly as much as the full sample; vendor metrics, age and body mass explained \rtwoRange\% of each component. Most of this variation arose within athletes, so landing information needs averaging over trials and sessions.

\medskip
\noindent\textbf{Keywords:} Countermovement jump; Principal component analysis; Force--time curve; Force plates; Athlete monitoring; Landing; Sports analytics; Football
